\subsubsection{Leseschlüssel}\label{leseschluxfcssel}

\begin{itemize}
\tightlist
\item
  \textbf{Exakt:} Identität oder mathematische Folgerung im angegebenen
  Modell.
\item
  \textbf{Numerisch:} berechneter Näherungswert, ohne zertifizierte
  Einschließung.
\item
  \textbf{Bedingt:} exakter Satz erst nach ausdrücklich zusätzlicher
  Ressource.
\item
  \textbf{Offen:} fehlende Herleitung oder physikalische Identifikation.
\end{itemize}

Eine grüne Prüfung ersetzt weder die Voraussetzungen eines Satzes noch
eine ausführbare Operation. Assoziative Algebradimension, dynamische
Kontrollierbarkeit und physische Verfügbarkeit werden getrennt geführt.

\subsection{1. Die unveränderte gemeinsame
Grundlage}\label{die-unveruxe4nderte-gemeinsame-grundlage}

Mit 64 Fermionmoden, 60 Bosonmoden und demselben gepinnten Tensor gilt

\[\adjustbox{max width=\linewidth}{$\displaystyle 
H_{\rm nat}=\Delta N_b+g\sum_A(b_A^\dagger P_A+P_A^\dagger b_A),
\qquad P_A=\sum_{i<j}W_{A,ij}f_jf_i,
\qquad N=N_f+2N_b.
$}\]

W besitzt 480 Einträge ±1, acht disjunkte Paare je Zeile, und
\codeword{WW†=8 I60}. Die Fermionmarkierung ist \codeword{16⊗4}, die
Bosonmarkierung \codeword{10⊗6}. Spin(10)×SU(4) ist zunächst eine
\textbf{innere} Quellsymmetrie, keine bereits hergeleitete
Raumzeit-Lorentzgruppe.

Der frühere native Satz bleibt unverändert maßgeblich: Für
Δ\textgreater0 und \codeword{0<|g|/Δ≤1/20} hat dieser μ=0-Vertrag einen
eindeutigen globalen Grundzustand Ω bei N=64. Am Prüfpunkt
\codeword{g/Δ=1/20} gelten unter anderem

\[\adjustbox{max width=\linewidth}{$\displaystyle 
-1.158089\Delta<E_0<-1.129636\Delta,
\qquad 0.842846<\langle N_b\rangle_\Omega<1.245656.
$}\]

Die isolierte niedrige Entnahmelinie im N=63-Sektor hat weiterhin die
bewiesenen Schranken aus v1.6.4/v1.6.5:

\[\adjustbox{max width=\linewidth}{$\displaystyle 
0.007737\Delta<\epsilon<0.039079764\Delta,
\qquad
Z_{\rm low}>\frac{40912436089}{46487375000}>0.88007628.
$}\]

Z bezieht sich auf das gesamte normierte Fermionspektralgewicht pro
Mode. Die gesamte Entnahmenorm ist dagegen
\codeword{ν=<f\_r†f\_r>=1-<Nb>/32}. Die beiden Größen sind nicht
austauschbar. Der Pol ist eine wohldefinierte isolierte Eigenlinie
dieses endlichen Modenmodells; seine relativistische Teilchen-, Massen-
oder räumliche Propagationsdeutung ist eine weitere Aufgabe.

Der vollständige bedingte Zwei-Banken-Satz aus v1.6.5 bleibt erhalten:
mit deklariertem Fermionlink und Rotorvertrag gelingt der reine niedrige
Transfer mit Wahrscheinlichkeit \textgreater99,267 \%, die unfiltrierte
ursprüngliche Entnahme mit \textgreater89,726 \%. Der Link war und ist
\textbf{zusätzlich vorausgesetzt}. Die neuen Vierzustandszahlen unten
gehören zu einem anderen Vertrag und ersetzen diese Ergebnisse nicht.

\subsection{2. Quellenprüfung der beiden neuen
Ergebnisrunden}\label{quellenpruxfcfung-der-beiden-neuen-ergebnisrunden}

\subsubsection{2.1 Bestätigte Zerlegung, korrigierte
Operationsaussagen}\label{bestuxe4tigte-zerlegung-korrigierte-operationsaussagen}

Die 944, 31 und 152 mitgelieferten Prüfbedingungen der Symmetrie-,
Verfügbarkeits- und Lorentzprogramme sind normal und optimiert
byteidentisch reproduziert. Darunter befinden sich neun wörtliche
\codeword{need(True, ...)}-Guards mit theoretischen Folgerungen. Sie
zählen als ausgeführte Guards, nicht als neun unabhängige Beweise. Die
Clock-Normalisatorbehauptung gehörte dazu; Abschnitt 3 ersetzt sie durch
eine explizite Konstruktion.

Im gesamten N=3-Raum gilt die Zerlegung:

{\def\LTcaptype{none} % do not increment counter
\begin{longtable}[]{@{}
  >{\raggedright\arraybackslash}p{(\linewidth - 8\tabcolsep) * \real{0.1579}}
  >{\raggedleft\arraybackslash}p{(\linewidth - 8\tabcolsep) * \real{0.2105}}
  >{\raggedleft\arraybackslash}p{(\linewidth - 8\tabcolsep) * \real{0.2105}}
  >{\raggedleft\arraybackslash}p{(\linewidth - 8\tabcolsep) * \real{0.2105}}
  >{\raggedleft\arraybackslash}p{(\linewidth - 8\tabcolsep) * \real{0.2105}}@{}}
\toprule\noalign{}
\begin{minipage}[b]{\linewidth}\raggedright
Typ
\end{minipage} & \begin{minipage}[b]{\linewidth}\raggedleft
irreduzible Dimension
\end{minipage} & \begin{minipage}[b]{\linewidth}\raggedleft
Multiplizität
\end{minipage} & \begin{minipage}[b]{\linewidth}\raggedleft
Spin-Casimir
\end{minipage} & \begin{minipage}[b]{\linewidth}\raggedleft
Farb-Casimir
\end{minipage} \\
\midrule\noalign{}
\endhead
\bottomrule\noalign{}
\endlastfoot
(1200,20) & 24000 & 1 & 141/4 & 39/4 \\
(560,20′) & 11200 & 1 & \textbf{117/4} & 63/4 \\
(672,4̄) & 2688 & 1 & 165/4 & 15/4 \\
(144,20) & 2880 & 2 & 85/4 & 39/4 \\
(144,4̄) & 576 & 2 & 85/4 & 15/4 \\
(16̄,20) & 320 & 2 & 45/4 & 39/4 \\
(16̄,4̄) & 64 & 1 & 45/4 & 15/4 \\
\end{longtable}
}

Die gewichtete Gesamtdimension ist 45504. Nur die getrennten Teile Λ³64
und 60⊗64 sind jeweils multiplizitätsfrei, nicht ihre direkte Summe. Die
drei dunklen Typen besitzen alle Gesamtcasimir 45, werden aber bereits
durch \textbf{einen} getrennten Casimir unterschieden.

Eine unabhängige rationale Rechnung auf den kleinen Multiplizitätsräumen
liefert für die jeweils \textbf{gewährten} Operationen:

{\def\LTcaptype{none} % do not increment counter
\begin{longtable}[]{@{}lr@{}}
\toprule\noalign{}
Operationsvertrag & komplexe assoziative Algebradimension \\
\midrule\noalign{}
\endhead
\bottomrule\noalign{}
\endlastfoot
Ein fixes H=Nb+X/20 & 8 \\
X und Nb unabhängig & 14 \\
Zusätzlich ein dunkler Isotypieprojektor & 15 \\
Zusätzlich ein getrennter Casimir & 16 \\
Zusätzlich nur Gesamtcasimir & 14 \\
\end{longtable}
}

Somit gibt es keine universelle binäre Frage „alle Symmetriegriffe oder
gar kein Fortschritt``. Die Dimension 15 ist ein konkretes
Gegenbeispiel. Die vollständige Gruppen-Generatoralgebra hat mit X,Nb
die Dimension 743583744 und Kommutantdimension 7. Sie ist unter
Gruppenkonjugation \textbf{stabil}, aber nicht punktweise invariant.
Sieben ist kein universeller Boden für beliebige symmetriestabile
Algebren: End(H) wäre ebenfalls stabil und hätte einen eindimensionalen
Kommutanten. Keine dieser Zahlen beweist, dass die zugehörigen
Operationen im Compiler ausführbar sind.

\subsubsection{2.2 Der Grundzustandsbericht enthält einen
Versionskonflikt}\label{der-grundzustandsbericht-enthuxe4lt-einen-versionskonflikt}

Der neue Text hat die Nichtschlussnorm bereits korrigiert, das
danebenliegende \codeword{native\_ground\_state.json} stammt jedoch von
einer älteren Codefassung. Sein als PASS markierter Wert ist um
\codeword{229²} zu groß; die dortige relative orthogonale Norm ist sogar
größer als eins. Der daraus berechnete Ritzwert −1,18826198 wird
\textbf{nicht} übernommen.

Eine neue direkte ganzzahlige Kontraktion umgeht die 15,25 Millionen
Komponenten von v₃. Sie berechnet \codeword{u=T T†v₂} auf 293280
nichtverschwindenden Komponenten und bestätigt exakt

\[\adjustbox{max width=\linewidth}{$\displaystyle 
\|u\|^2=752194252800,\quad
\langle v_2,u\rangle=575078400,\quad
w_2=u-\frac{299520}{229}v_2,
\quad\|w_2\|^2=\frac{5001523200}{229}>0.
$}\]

Der eine Vektor je Bosonlage schließt daher nicht. w₂ ist ein zweiter
Singulettvektor derselben Lage. Für die kleine vier- bzw.
fünfdimensionale Ritzmatrix ergeben sich bei g/Δ=1/20 \textbf{numerisch}

\[\adjustbox{max width=\linewidth}{$\displaystyle 
E_{R,4}/\Delta=-1.0942308394409612,
\qquad E_{R,5}/\Delta=-1.0942318710142458.
$}\]

Die zusätzliche w₂-Richtung ändert diese Näherung tatsächlich erst um
etwa 10⁻⁶. Das bedeutet \textbf{keine} Konvergenz zum wahren
Grundzustand: Beide Werte liegen noch über dessen bereits bewiesener
Obergrenze.

Die neue einfache Restformel zeigt die Ursache. Für den normierten
Vier-Vektor-Ritzeigenzustand mit letzter Komponente c₃ gilt

\[\adjustbox{max width=\linewidth}{$\displaystyle 
\|(H-E_R)\psi_3\|^2
=g^2|c_3|^2\frac{\nu_4+\|w_2\|^2}{\nu_3}.
$}\]

Mit dem früher geprüften vierten Moment ist die Restnorm numerisch rund
0,373063 Δ. Der kleine w₂-Zweig trägt weniger als 23 Millionstel des
quadrierten Restes; die ausgelassene v₄-Richtung dominiert. Die neuen
Zahlen \codeword{Zh≈0,9737} und \codeword{εh≈0,031Δ} gehören ebenfalls
zur Ritz-Näherung: Gesamtentnahmenorm und erstes Moment, nicht neu
bestimmte Polgrößen auf Ω.

Die Suche nach \codeword{E<−1,2Δ} oder \codeword{E<−1,1625Δ} wäre sogar
mit der bestehenden unteren Schranke unvereinbar. Die schwächere neue
Beweismethode öffnet den bereits geschlossenen nativen Grundzustandssatz
nicht wieder.

\subsection{3. Neuer exakter Anschluss: Der Clock liegt bereits in
Spin(10)}\label{neuer-exakter-anschluss-der-clock-liegt-bereits-in-spin10}

\subsubsection{3.1 Ausgangspunkt ist der tatsächliche
Quell-Clock}\label{ausgangspunkt-ist-der-tatsuxe4chliche-quell-clock}

Der gepinnte endliche Compilerpräfix liefert auf fünf Hilfsachsen

\[\adjustbox{max width=\linewidth}{$\displaystyle 
p=(0\mapsto2,\ 1\mapsto0,\ 2\mapsto1,\ 3\mapsto4,\ 4\mapsto3),
\qquad\det p=-1.
$}\]

Seine vorzeichenrichtige Exteriorhebung auf dem geraden
5-Moden-Hilfsfockraum ist Γ(p)\textbar even. Auf den nativen Fermionen
ist \codeword{GF=Γ(p)|even⊗I4}; auf den Bosonen ist
\codeword{GB=(-P10)⊗I6}, wobei P10 beide Fünfergruppen permutiert. Diese
Darstellung wird aus dem vorhandenen Clockpräfix erneut extrahiert,
nicht aufgrund passender Eigenwerte ausgewählt.

\subsubsection{3.2 Ein ausdrückliches
Spinwort}\label{ein-ausdruxfcckliches-spinwort}

Auf dem 32-dimensionalen Hilfsfockraum seien

\[\adjustbox{max width=\linewidth}{$\displaystyle 
\gamma_j=a_j+a_j^\dagger,\qquad
\gamma_{j+5}=i(a_j^\dagger-a_j),\qquad 0\le j<5,
$}\]

und Π dessen Fermionparität. Setze

\[\adjustbox{max width=\linewidth}{$\displaystyle 
\Omega_{10}=\gamma_0\gamma_1\cdots\gamma_9=i\Pi,
\qquad
R_{ab}=\frac{(\gamma_a-\gamma_b)(\gamma_{a+5}-\gamma_{b+5})}{2}.
$}\]

R\_ab ist das Produkt zweier reeller Einheitsvektoren der
Cliffordalgebra, also ein Spin(10)-Element. Direkt gilt
\codeword{R\_ab=i Γ((ab))}. Da \codeword{p=(01)(02)(34)} mit rechts
zuerst angewandter Permutation, folgt

\[\adjustbox{max width=\linewidth}{$\displaystyle 
\boxed{S=\Omega_{10}R_{01}R_{02}R_{34}=\Pi\Gamma(p)\in\mathrm{Spin}(10).}
$}\]

S ist ein Produkt von 16 reellen Clifford-Einheitsvektoren. Die
Matrixprüfung bestätigt ohne Rundungstoleranzen

\[\adjustbox{max width=\linewidth}{$\displaystyle 
S_{even}=\Gamma(p)_{even},\qquad S_{odd}=-\Gamma(p)_{odd},
\qquad S\gamma_jS^\dagger=-\gamma_{p(j)}
$}\]

mit entsprechender Fortsetzung auf die zweite Fünfergruppe. Damit
stimmen \textbf{beide nativen Darstellungen} und die ursprüngliche
Kopplung überein:

\[\adjustbox{max width=\linewidth}{$\displaystyle 
G_F=S_{even}\otimes I_4,\qquad
G_B=(-P_{10})\otimes I_6,\qquad
W\Lambda^2G_F=G_BW.
$}\]

S hat Ordnung sechs. Die Konjugation aller 45 Spin-Erzeuger und die
Kommutation mit allen 15 Farb-Erzeugern wurden zusätzlich direkt
geprüft. Die Aussage ist stärker und präziser als ein lediglich
behaupteter Normalisatorsatz.

\textbf{Bedeutung:} Der dokumentierte Clock und die innere
Spin-Symmetrie brauchen hier keine getrennten zusätzlichen
mathematischen Träger. Es handelt sich um einen konkreten endlichen
Schritt derselben Darstellung. Daraus folgen aber weder kontinuierliche
Spin-Kontrollen noch die Identifikation des Clocks mit Hamiltonzeit. Die
Hilfs-Cliffordachsen sind nicht bereits fünf Raumrichtungen.

\subsection{4. Clock und getrennte Casimire lassen sich gemeinsam
auslesen}\label{clock-und-getrennte-casimire-lassen-sich-gemeinsam-auslesen}

Weil G aus Spin(10) stammt und auf SU(4) trivial wirkt, kommutiert es
mit beiden getrennten Casimiren, auch auf ihren Fockhebungen. Die Frage
aus dem neuen Text ist damit auf algebraischer Ebene beantwortet.

Die exakten Charakterfolgen von S\_even, S\_odd und dem Vektor sind

\[\adjustbox{max width=\linewidth}{$\displaystyle 
(16,0,4,0,4,0),\quad(16,0,4,0,4,0),\quad(10,0,4,-6,4,0).
$}\]

Die schon unabhängig nachgerechneten Zerlegungen
\codeword{Sym³16=672+144}, \codeword{Λ³16=560},
\codeword{S21(16)=1200+144+16̄} und \codeword{10⊗16=144+16̄} bestimmen die
Charaktere der sieben Typen. Die diskrete Fourierinversion erfolgt exakt
in Z{[}ζ6{]}, nicht durch gerundete komplexe Eigenwerte. Für die drei
zuvor gemeinsam dunklen Räume ergibt sich:

{\def\LTcaptype{none} % do not increment counter
\begin{longtable}[]{@{}
  >{\raggedright\arraybackslash}p{(\linewidth - 12\tabcolsep) * \real{0.1111}}
  >{\raggedleft\arraybackslash}p{(\linewidth - 12\tabcolsep) * \real{0.1481}}
  >{\raggedleft\arraybackslash}p{(\linewidth - 12\tabcolsep) * \real{0.1481}}
  >{\raggedleft\arraybackslash}p{(\linewidth - 12\tabcolsep) * \real{0.1481}}
  >{\raggedleft\arraybackslash}p{(\linewidth - 12\tabcolsep) * \real{0.1481}}
  >{\raggedleft\arraybackslash}p{(\linewidth - 12\tabcolsep) * \real{0.1481}}
  >{\raggedleft\arraybackslash}p{(\linewidth - 12\tabcolsep) * \real{0.1481}}@{}}
\toprule\noalign{}
\begin{minipage}[b]{\linewidth}\raggedright
Dunkler Typ
\end{minipage} & \begin{minipage}[b]{\linewidth}\raggedleft
Phase 0
\end{minipage} & \begin{minipage}[b]{\linewidth}\raggedleft
1
\end{minipage} & \begin{minipage}[b]{\linewidth}\raggedleft
2
\end{minipage} & \begin{minipage}[b]{\linewidth}\raggedleft
3
\end{minipage} & \begin{minipage}[b]{\linewidth}\raggedleft
4
\end{minipage} & \begin{minipage}[b]{\linewidth}\raggedleft
5
\end{minipage} \\
\midrule\noalign{}
\endhead
\bottomrule\noalign{}
\endlastfoot
(1200,20) & 4000 & 4000 & 4000 & 4000 & 4000 & 4000 \\
(560,20′) & 1920 & 1840 & 1840 & 1920 & 1840 & 1840 \\
(672,4̄) & 464 & 440 & 440 & 464 & 440 & 440 \\
Summe, bisherige Clock-Auslesung & 6384 & 6280 & 6280 & 6384 & 6280 &
6280 \\
\end{longtable}
}

Ein separater Casimir trennt jede der sechs dunklen Clockzellen in genau
drei Teile. Die hellen Zweier-Multiplizitätsräume bleiben erhalten.
Daraus folgt im N=3-Sektor

\[\adjustbox{max width=\linewidth}{$\displaystyle 
\boxed{\dim\operatorname{Alg}^*(X,N_b,G,C_S)=96,}
\qquad
\boxed{\dim\operatorname{Alg}^*(X,N_b,G,C_S)'=119597824.}
$}\]

Zum Vergleich: ohne C\_S sind es 84 und 240742144. C\_C kann C\_S
ersetzen. Die Lagrangeprojektoren auf den drei dunklen C\_S-Werten
bestätigen die Trennung ausdrücklich; beide Casimire zugleich werden
nicht benötigt.

\textbf{Es handelt sich um eine gelöste Kombinationsfrage, nicht um eine
neue Kontrollfreigabe.} Der benötigte Casimir ist noch kein vom Compiler
bereitgestelltes Messinstrument. Der große verbleibende Kommutant zeigt
weiterhin viele nicht aufgelöste Freiheitsgrade. Die Dimensionszahlen
gelten für N=3, nicht als vollständige Klassifikation aller
Ladungssektoren.

